\documentclass[12pt, a4paper]{article}

\usepackage[a4paper, top = 1 in, bottom = 1 in, left = 1 in, right = 1 in]{geometry}

\usepackage{amsmath, amssymb, amsfonts}
\allowdisplaybreaks[4]

\usepackage{enumerate}
\usepackage{multirow}
\usepackage{hhline}
\usepackage{array}
\usepackage{longtable}
\usepackage{graphicx}
\usepackage{tabularray}
\usepackage{undertilde}
\usepackage{xcolor}
\usepackage{dingbat}
\usepackage{fontawesome5}
\usepackage[colorlinks=true, linkcolor=blue, urlcolor=red]{hyperref}
\usepackage{tasks}
\usepackage{bbding}
\usepackage{twemojis}
% how to use bull's eye ----- \scalebox{2.0}{\twemoji{bullseye}}
\usepackage{customdice}
% how to put dice face ------ \dice{2}
\usepackage{fontspec}
\usepackage{fancyhdr}
\usepackage{lastpage}



\pagestyle{fancy}
\fancyhf{}
\fancyfoot[C]{Page \thepage\ of \pageref{LastPage}}
\renewcommand{\headrulewidth}{0pt}



\title{MSMS 401 : Bayesian Statistics and Computation}
\author{{\fontsize{25}{20} \benyorithafont Ananda Biswas}}
\date{}

\newfontface\gloriahallelujahfont{GLORIAHALLELUJAH-REGULAR.TTF}

\newfontface\benyorithafont{Benyoritha-G334D.otf}


\begin{document}

\maketitle

\section*{\faArrowAltCircleRight[regular] \textcolor{blue}{Problem}}

Consider a random sample of size $n$ from $N(\mu_0, \sigma^2)$, $\mu_0^2$ is the known mean. Take a non-informative prior for $\lambda = \sigma^2$. Obtain the posterior distribution of $\lambda$, hence obtain Bayes estimator of $\lambda$ assuming squared-error loss function.

\section*{\faArrowAltCircleRight[regular] \textcolor{blue}{Solution}}

Let $x_1, x_2, \ldots, x_n \overset{iid}{\sim} f(x \mid \lambda)$ where

\begin{align*}
f(x \mid \lambda) = \dfrac{1}{\sqrt{\lambda} \cdot \sqrt{2\pi}} \exp\left\{ -\dfrac{(x-\mu_0)^2}{2\lambda} \right\}.
\end{align*}

Jeffreys prior for $\lambda$ will be $\lambda \sim \pi(\lambda) = \sqrt{I(\lambda)}$. Now, 

\begin{align*}
f(x \mid \lambda) &= \dfrac{1}{\sqrt{\lambda} \cdot \sqrt{2\pi}} \exp\left\{ -\dfrac{(x-\mu_0)^2}{2\lambda} \right\} \\[1em]
\Rightarrow \ln f(x \mid \lambda) &= - \dfrac{1}{2} \ln \lambda + \ln \dfrac{1}{\sqrt{2\pi}} -\dfrac{(x-\mu_0)^2}{2\lambda} \\[1em]
\Rightarrow \dfrac{\partial}{\partial \lambda} \ln f(x \mid \lambda) &= - \dfrac{1}{2\lambda} + \dfrac{(x-\mu_0)^2}{2\lambda^2} \\[1em]
\Rightarrow \dfrac{\partial^2}{\partial \lambda^2} \ln f(x \mid \lambda)  &= \dfrac{1}{2\lambda^2} - \dfrac{(x-\mu_0)^2}{\lambda^3}.
\end{align*}

\begin{align*}
\therefore I(\lambda) &= -E\left[ \dfrac{\partial^2}{\partial \lambda^2} \ln f(X \mid \lambda) \right] \\[1em]
&= - \dfrac{1}{2\lambda^2} + \dfrac{E(X - \mu_0)^2}{\lambda^2} \\[1em]
&= - \dfrac{1}{2\lambda^2} + \dfrac{1}{\lambda^2} \\[1em]
&= \dfrac{1}{2\lambda^2}
\end{align*}

\begin{align*}
\therefore \pi(\lambda) = \sqrt{\dfrac{1}{2\lambda^2}} = \dfrac{1}{\sqrt{2} \cdot \lambda}, \; \lambda > 0.
\end{align*}

This is improper, so we can ignore any multiplicative constants and take

\begin{align*}
\pi(\lambda) \propto \dfrac{1}{\lambda}.
\end{align*}



The likelihood for $\lambda$ is

\begin{align*}
f(\mathbf{x} \mid \lambda) &= \prod \limits_{i = 1}^{n} \dfrac{1}{\sqrt{\lambda} \cdot \sqrt{2\pi}} \exp\left\{ -\dfrac{(x_i-\mu_0)^2}{2\lambda} \right\} \\[1em]
&= \left( \dfrac{1}{\sqrt{\lambda} \cdot \sqrt{2\pi}} \right)^n \cdot \exp\left\{ -\dfrac{1}{2\lambda} \sum \limits_{i = 1}^{n}(x_i-\mu_0)^2 \right\}
\end{align*}

We compute the posterior distribution of $\lambda$ as
\begin{align*}
p(\lambda \mid \mathbf{x}) &= \dfrac{f(\mathbf{x} \mid \lambda) \, \cdot \, \pi(\lambda)}{\displaystyle \int\limits_{0}^{\infty} f(\mathbf{x} \mid \lambda) \, \cdot \, \pi(\lambda) \, d\lambda} \\[1em]
&= \dfrac{\dfrac{1}{\lambda} \cdot \left( \dfrac{1}{\sqrt{\lambda} \cdot \sqrt{2\pi}} \right)^n \cdot \exp\left\{ -\dfrac{1}{2\lambda} \sum \limits_{i = 1}^{n}(x_i-\mu_0)^2 \right\}}{\displaystyle \int\limits_{0}^{\infty} \dfrac{1}{\lambda} \cdot \left( \dfrac{1}{\sqrt{\lambda} \cdot \sqrt{2\pi}} \right)^n \cdot \exp\left\{ -\dfrac{1}{2\lambda} \sum \limits_{i = 1}^{n}(x_i-\mu_0)^2 \right\} \, d\lambda}.
\end{align*}

Denote $S := \sum \limits_{i = 1}^{n}(x_i-\mu_0)^2$. Then

\begin{align*}
p(\lambda \mid \mathbf{x}) = \dfrac{\lambda^{-\left(\dfrac{n}{2} + 1\right)} \cdot \exp\left\{- \dfrac{S/2}{\lambda} \right\}}{\displaystyle \int\limits_{0}^{\infty} \lambda^{-\left(\dfrac{n}{2} + 1\right)} \cdot \exp\left\{- \dfrac{S/2}{\lambda} \right\} \, d\lambda}.
\end{align*}

\newpage

Using the result $\displaystyle \int\limits_{0}^{\infty} x^{-(\alpha + 1)} \cdot \exp\left\{- \dfrac{\beta}{x} \right\} \, dx = \dfrac{\Gamma(\alpha)}{\beta^{\alpha}}$, we get

\begin{align*}
p(\lambda \mid \mathbf{x}) = \dfrac{\left(\dfrac{S}{2}\right)^{\dfrac{n}{2}}}{\Gamma\left(\dfrac{n}{2}\right)} \cdot \lambda^{-\left(\dfrac{n}{2} + 1\right)} \cdot \exp\left\{- \dfrac{S/2}{\lambda} \right\}
\end{align*}

which we identify as the PDF of Inverse-Gamma$\left(\dfrac{n}{2}, \dfrac{S}{2} \right)$.

\vspace{0.5cm}

$\therefore$ \textit{a posteriori} $\lambda \mid \mathbf{x} \sim \text{Inverse-Gamma}\left(\dfrac{n}{2}, \dfrac{S}{2} \right)$.

\vspace{0.5cm}

{\fontsize{12}{20} \gloriahallelujahfont Under squared error loss function, Bayes estimate is the posterior mean.}

\vspace{0.2cm}

The expected value of an Inverse-Gamma$(\alpha, \beta)$ variate is $\dfrac{\beta}{\alpha - 1}$. Thus,

\begin{align*}
\hat{\lambda}_{Bayes} = \dfrac{S}{n-1} = \dfrac{\sum \limits_{i = 1}^{n}(x_i-\mu_0)^2}{n-1}.
\end{align*}
\end{document}